\documentclass[11pt]{article}
\usepackage[margin=1in]{geometry}
\usepackage{amsmath,amssymb}
\setlength{\parindent}{0pt}
\begin{document}
\section*{HW08 Question 7}

\textbf{Problem.} FNE scores are compared for bulimic and normal female students. Technology output gives difference \(=\mu(\text{Bulimic})-\mu(\text{Normal})\), estimate \(3.68\), 95\% CI \((-0.60,7.95)\), \(T=1.79\), \(p=0.086\), \(df=23\).

\textbf{Solution.}

The requested 95\% confidence interval is read directly from the technology output:
\[
\boxed{(-0.60,\ 7.95)}.
\]
For the assumptions part, select A, C, and D:
\[
\boxed{\text{A, C, and D}}
\]
A: the scores were randomly selected in an independent manner. C: both group distributions are approximately normal. D: the two population variances are equal. Do not select B because \(n_1=11\) and \(n_2=14\), so both samples are not large.

For the conclusion, compare the interval and P-value:
\[
0\in(-0.60,7.95),\qquad p=0.086>0.05.
\]

\[
\boxed{\text{Assumptions: A, C, and D. The assumptions are reasonably satisfied.}}
\]
The interval contains 0 and the P-value is greater than 0.05, so there is not sufficient evidence of a difference in mean FNE score.
\end{document}
